\section{Evaluating the Integral in Detail}
\label{app:integral}

\paragraph{Trapezoid rule.}
We approximate the integral of $g$ from \eqref{eq:cfi} with step $h$ and $K$ nodes $s_k = kh$. Every node gets the weight $h$, except $s_0 = 0$, which gets $h/2$.
For small $s$ we have $\mathrm{Im}\, \phi_z(s) \approx s\, \E[z \mid \xo]$ and $\sinh(\pi s) \approx \pi s$, so $g(0) = \mu_z / \pi$, which gives the term $\frac{\mu_z}{2\pi}$ in \eqref{eq:trapezoid}.
The integrand is smooth, so the trapezoid rule converges exponentially fast in the number of nodes.
Two things have to be chosen: where to stop, $S = Kh$, and the step $h$.

\paragraph{Truncation.}
Since $|\phi_z| \leq 1$, the part of the integral beyond $S$ is bounded independently of the density:
\[
  \int_S^\infty \frac{\mathrm{d}s}{\sinh(\pi s)} = \frac{1}{\pi} \log \coth \frac{\pi S}{2} \approx \frac{2}{\pi} e^{-\pi S}.
\]
For $S = 12$ this is about $3 \cdot 10^{-17}$, which is below double precision, so we never need frequencies beyond $S = 12$.

\paragraph{Step size.}
For a fixed $z$, the trapezoid sum $h \sum_k \frac{\sin(khz)}{\sinh(\pi kh)}$ is periodic in $z$ with period $\frac{2\pi}{h}$.
So it can only approximate the sigmoid on one period, roughly for $|z| < \frac{\pi}{h}$, and outside of it the sum repeats instead of staying close to $0$ or $1$.
The expected prediction is therefore accurate as long as $p(z \mid \xo)$ has nearly all its mass within $|z| < \frac{\pi}{h}$.
In \eqref{eq:step}, the period covers $\mu_z \pm 6\sigma_z$, with a margin of $1$ for the width of the sigmoid itself, and the nodes reach $S$.
If the first term is the smaller one, the nodes stop before $S$. For a wide $p(z \mid \xo)$ this does not matter, since the CF then decays quickly. For a narrow $p(z \mid \xo)$ far away from zero, $K$ has to be large enough:
\[
  K \geq S (|\mu_z| + 6\sigma_z + 1) / \pi.
\]
The step only affects the error and not the value that is approximated, so it can be chosen for every point separately, and it needs no gradient in training.
The mean and variance of $z$ follow from the mean and covariance of the missing features,
\[
  \mu_z = w_o^T \xo + w_m^T \E[\xm \mid \xo], \qquad \sigma_z^2 = w_m^T\, \mathrm{Cov}[\xm \mid \xo]\, w_m,
\]
which are tractable for GMMs and PCs (\cref{app:density}).

\paragraph{Cost.}
The method needs $K$ evaluations of the CF per point. All of them are at multiples $s_k w_m$ of the same vector, so they can be done in one batched pass through the density, plus one pass for the moments. In our experiments $K = 40$ is enough. % TODO check against the frequency plot

\section{Characteristic Function and Moments of the Density}
\label{app:density}

\paragraph{Probabilistic circuits.}
A PC is a computational graph of sum nodes, product nodes and leaves, where every leaf is a simple distribution over a single feature \citep{pc_intro, spn}.
If the circuit is smooth and decomposable, the expectation of any function that factorises over the features, $f(x) = \prod_d f_d(x_d)$, is computed in one pass, with every leaf replaced by $\E_{\text{leaf}}[f_d(x_d)]$.
A GMM with diagonal covariances is the simplest such circuit, a single sum node over products of Gaussian leaves.

\paragraph{The conditional CF of a PC.}
The CF is the expectation of $e^{i s^T x} = \prod_d e^{i s_d x_d}$, which factorises over the features, so a PC computes it in one pass \citep{yu2023characteristic}:
\[
  \phi_{\xm \mid \xo}(s) = \frac{1}{p(\xo)} \int e^{i s^T \xm}\, p(\xo, \xm)\, \mathrm{d}\xm.
\]
Every leaf of an observed feature returns its density at $x_d$, and every leaf of a missing feature returns its own CF at $s_d$. For a Gaussian leaf this is $\exp(i s_d \mu_d - \frac{1}{2} s_d^2 \sigma_d^2)$, and for a categorical leaf $\sum_v p(x_d{=}v)\, e^{i s_d v}$.
Evaluating the circuit bottom-up gives $\int e^{i s^T \xm} p(\xo, \xm)\, \mathrm{d}\xm = p(\xo)\, \phi_{\xm \mid \xo}(s)$, and dividing by $p(\xo)$ gives the conditional CF.
The values are complex, and products over many leaves can underflow, so we evaluate the circuit in log space, where products become sums and sums become a complex log-sum-exp.

\paragraph{The moments of $z$ in a PC.}
The mean and variance of $z$ are computed with the same kind of pass.
Every node $n$ carries three numbers: its value $P_n$, as in the marginal likelihood, and the first and second moments $M_n = \int z_n\, p_n$ and $Q_n = \int z_n^2\, p_n$ of the part $z_n$ of $z$ over its scope.
A sum node takes the weighted sums of the three numbers of its children.
For a product node with children $a$ and $b$, we have $z_n = z_a + z_b$ with independent parts, so
\[
  P_n = P_a P_b, \qquad M_n = M_a P_b + P_a M_b, \qquad Q_n = Q_a P_b + 2 M_a M_b + P_a Q_b.
\]
A leaf of an observed feature has $P = p(x_d)$, $M = w_d x_d\, p(x_d)$ and $Q = w_d^2 x_d^2\, p(x_d)$, and a Gaussian leaf of a missing feature has $P = 1$, $M = w_d \mu_d$ and $Q = w_d^2 (\sigma_d^2 + \mu_d^2)$.
At the root, $\mu_z = M / P$ and $\sigma_z^2 = Q / P - \mu_z^2$.

\paragraph{Diagonal GMM.}
For a GMM with diagonal covariances, the pass reduces to the closed form
\[
  \phi_{\xm \mid \xo}(s) = \sum_{k=1}^{K} \tilde{\pi}_k(\xo)\, e^{i s^T \mu_k^m - \frac{1}{2} s^T \Sigma_k^m s},
\]
with the responsibilities $\tilde{\pi}_k(\xo) \propto \pi_k\, \mathcal{N}(\xo; \mu_k^o, \Sigma_k^o)$, and the means and covariances $\mu_k^m, \Sigma_k^m$ of the components restricted to the missing features.
The responsibilities only depend on $\xo$, so they are computed once per input and reused for every frequency.
Within component $k$, $z$ is Gaussian, $z \sim \mathcal{N}(m_k, v_k)$, with $m_k = w_o^T \xo + w_m^T \mu_k^m$ and $v_k = w_m^T \Sigma_k^m w_m$. So
\[
  \mu_z = \sum_k \tilde{\pi}_k(\xo)\, m_k, \qquad \sigma_z^2 = \sum_k \tilde{\pi}_k(\xo) \left( v_k + (m_k - \mu_z)^2 \right),
\]
where the variance is the variance within the components plus the spread of their means.

\section{Exact Reference for a GMM}
\label{app:gmm-reference}

In the experiments we compare to the exact expected prediction under a diagonal GMM, which can be computed without the CF.
Within component $k$, $z \sim \mathcal{N}(m_k, v_k)$ (\cref{app:density}), so $p(z \mid \xo)$ is a mixture of Gaussians and
\[
  p(y{=}1 \mid \xo) = \sum_{k=1}^{K} \tilde{\pi}_k(\xo)\, \E_{z \sim \mathcal{N}(m_k, v_k)} \left[ \sigma(z) \right].
\]
Each term is a one-dimensional logistic-normal integral, which we compute to machine precision with adaptive quadrature (\texttt{scipy.integrate.quad}).
This does not work for general PCs: the number of mixture components of a PC grows exponentially with its size, and with non-Gaussian leaves $z$ is not Gaussian within a component.

\section{Further Results on Real Data}
\label{app:results}

The setup is the same as in \cref{sec:results-real}, with 10 and 50\,\% missing.

\begin{table}[h]
  \centering
  \caption{Log-loss on the incomplete test points, 10\,\% missing.}
  \label{tab:real-log-loss-10}
  \resizebox{0.9\textwidth}{!}{\begin{tabular}{lcccccc}
\toprule
Dataset & No missing & Mean imp. & Cond.\ mean imp. & MICE & Generative & PCLR \\
\midrule
magic & 0.483{\scriptsize$\pm$0.005} & 0.509{\scriptsize$\pm$0.004} & 0.498{\scriptsize$\pm$0.003} & 0.500{\scriptsize$\pm$0.003} & 0.512{\scriptsize$\pm$0.005} & \textbf{0.497{\scriptsize$\pm$0.003}} \\
phoneme & 0.584{\scriptsize$\pm$0.007} & 0.594{\scriptsize$\pm$0.008} & 0.590{\scriptsize$\pm$0.008} & 0.594{\scriptsize$\pm$0.008} & \textbf{0.446{\scriptsize$\pm$0.008}} & 0.589{\scriptsize$\pm$0.009} \\
spambase & 0.258{\scriptsize$\pm$0.018} & \textbf{0.266{\scriptsize$\pm$0.024}} & 0.268{\scriptsize$\pm$0.022} & 0.272{\scriptsize$\pm$0.026} & 0.402{\scriptsize$\pm$0.028} & 0.271{\scriptsize$\pm$0.024} \\
banknote & 0.071{\scriptsize$\pm$0.009} & 0.180{\scriptsize$\pm$0.014} & \textbf{0.135{\scriptsize$\pm$0.011}} & 0.136{\scriptsize$\pm$0.005} & 0.334{\scriptsize$\pm$0.071} & 0.143{\scriptsize$\pm$0.013} \\
qsar & 0.474{\scriptsize$\pm$0.076} & 0.528{\scriptsize$\pm$0.067} & 0.516{\scriptsize$\pm$0.065} & \textbf{0.492{\scriptsize$\pm$0.065}} & 0.570{\scriptsize$\pm$0.029} & 0.528{\scriptsize$\pm$0.085} \\
pima & 0.563{\scriptsize$\pm$0.058} & 0.572{\scriptsize$\pm$0.054} & 0.567{\scriptsize$\pm$0.052} & 0.571{\scriptsize$\pm$0.059} & 0.590{\scriptsize$\pm$0.009} & \textbf{0.566{\scriptsize$\pm$0.052}} \\
\bottomrule
\end{tabular}
}
\end{table}

\begin{table}[h]
  \centering
  \caption{Log-loss on the incomplete test points, 50\,\% missing.}
  \label{tab:real-log-loss-50}
  \resizebox{0.9\textwidth}{!}{\begin{tabular}{lcccccc}
\toprule
Dataset & No missing & Mean imp. & Cond.\ mean imp. & MICE & Generative & PCLR \\
\midrule
magic & 0.483{\scriptsize$\pm$0.005} & 0.592{\scriptsize$\pm$0.008} & 0.556{\scriptsize$\pm$0.005} & 0.577{\scriptsize$\pm$0.004} & \textbf{0.541{\scriptsize$\pm$0.003}} & 0.544{\scriptsize$\pm$0.007} \\
phoneme & 0.584{\scriptsize$\pm$0.007} & 0.631{\scriptsize$\pm$0.009} & 0.617{\scriptsize$\pm$0.009} & 0.634{\scriptsize$\pm$0.010} & \textbf{0.499{\scriptsize$\pm$0.009}} & 0.603{\scriptsize$\pm$0.010} \\
spambase & 0.258{\scriptsize$\pm$0.018} & \textbf{0.352{\scriptsize$\pm$0.022}} & 0.378{\scriptsize$\pm$0.012} & 0.367{\scriptsize$\pm$0.024} & 0.447{\scriptsize$\pm$0.028} & 0.355{\scriptsize$\pm$0.019} \\
banknote & 0.071{\scriptsize$\pm$0.009} & 0.452{\scriptsize$\pm$0.022} & 0.419{\scriptsize$\pm$0.046} & 0.453{\scriptsize$\pm$0.041} & 0.445{\scriptsize$\pm$0.038} & \textbf{0.413{\scriptsize$\pm$0.030}} \\
qsar & 0.474{\scriptsize$\pm$0.076} & 0.595{\scriptsize$\pm$0.057} & 0.567{\scriptsize$\pm$0.062} & \textbf{0.528{\scriptsize$\pm$0.047}} & 0.544{\scriptsize$\pm$0.041} & 0.541{\scriptsize$\pm$0.048} \\
pima & 0.563{\scriptsize$\pm$0.058} & 0.612{\scriptsize$\pm$0.035} & 0.618{\scriptsize$\pm$0.037} & 0.611{\scriptsize$\pm$0.036} & 0.607{\scriptsize$\pm$0.011} & \textbf{0.606{\scriptsize$\pm$0.042}} \\
\bottomrule
\end{tabular}
}
\end{table}

\begin{table}[h]
  \centering
  \caption{Accuracy on the incomplete test points, 10\,\% missing.}
  \label{tab:real-accuracy-10}
  \resizebox{0.9\textwidth}{!}{\begin{tabular}{lcccccc}
\toprule
Dataset & No missing & Mean imp. & Cond.\ mean imp. & MICE & Generative & PCLR \\
\midrule
magic & 0.782{\scriptsize$\pm$0.005} & 0.763{\scriptsize$\pm$0.004} & 0.771{\scriptsize$\pm$0.002} & 0.769{\scriptsize$\pm$0.005} & 0.758{\scriptsize$\pm$0.003} & \textbf{0.771{\scriptsize$\pm$0.002}} \\
phoneme & 0.721{\scriptsize$\pm$0.009} & 0.713{\scriptsize$\pm$0.010} & 0.718{\scriptsize$\pm$0.009} & 0.715{\scriptsize$\pm$0.008} & \textbf{0.750{\scriptsize$\pm$0.006}} & 0.720{\scriptsize$\pm$0.008} \\
spambase & 0.921{\scriptsize$\pm$0.009} & 0.914{\scriptsize$\pm$0.009} & 0.912{\scriptsize$\pm$0.009} & 0.910{\scriptsize$\pm$0.010} & 0.840{\scriptsize$\pm$0.027} & \textbf{0.918{\scriptsize$\pm$0.010}} \\
banknote & 0.976{\scriptsize$\pm$0.004} & 0.932{\scriptsize$\pm$0.007} & 0.944{\scriptsize$\pm$0.009} & \textbf{0.945{\scriptsize$\pm$0.006}} & 0.863{\scriptsize$\pm$0.041} & 0.943{\scriptsize$\pm$0.010} \\
qsar & 0.814{\scriptsize$\pm$0.015} & 0.782{\scriptsize$\pm$0.020} & 0.786{\scriptsize$\pm$0.016} & \textbf{0.816{\scriptsize$\pm$0.013}} & 0.700{\scriptsize$\pm$0.027} & 0.783{\scriptsize$\pm$0.014} \\
pima & 0.734{\scriptsize$\pm$0.053} & 0.729{\scriptsize$\pm$0.041} & \textbf{0.731{\scriptsize$\pm$0.046}} & 0.725{\scriptsize$\pm$0.057} & 0.661{\scriptsize$\pm$0.021} & 0.727{\scriptsize$\pm$0.049} \\
\bottomrule
\end{tabular}
}
\end{table}

\begin{table}[h]
  \centering
  \caption{Accuracy on the incomplete test points, 50\,\% missing.}
  \label{tab:real-accuracy-50}
  \resizebox{0.9\textwidth}{!}{\begin{tabular}{lcccccc}
\toprule
Dataset & No missing & Mean imp. & Cond.\ mean imp. & MICE & Generative & PCLR \\
\midrule
magic & 0.782{\scriptsize$\pm$0.005} & 0.699{\scriptsize$\pm$0.008} & 0.741{\scriptsize$\pm$0.008} & 0.715{\scriptsize$\pm$0.008} & 0.742{\scriptsize$\pm$0.002} & \textbf{0.745{\scriptsize$\pm$0.007}} \\
phoneme & 0.721{\scriptsize$\pm$0.009} & 0.654{\scriptsize$\pm$0.015} & 0.679{\scriptsize$\pm$0.015} & 0.661{\scriptsize$\pm$0.018} & \textbf{0.730{\scriptsize$\pm$0.009}} & 0.700{\scriptsize$\pm$0.011} \\
spambase & 0.921{\scriptsize$\pm$0.009} & \textbf{0.870{\scriptsize$\pm$0.014}} & 0.848{\scriptsize$\pm$0.011} & 0.869{\scriptsize$\pm$0.011} & 0.820{\scriptsize$\pm$0.017} & 0.850{\scriptsize$\pm$0.010} \\
banknote & 0.976{\scriptsize$\pm$0.004} & 0.737{\scriptsize$\pm$0.017} & \textbf{0.781{\scriptsize$\pm$0.025}} & 0.753{\scriptsize$\pm$0.023} & 0.777{\scriptsize$\pm$0.027} & 0.778{\scriptsize$\pm$0.013} \\
qsar & 0.814{\scriptsize$\pm$0.015} & 0.743{\scriptsize$\pm$0.009} & 0.762{\scriptsize$\pm$0.016} & 0.763{\scriptsize$\pm$0.023} & 0.696{\scriptsize$\pm$0.048} & \textbf{0.768{\scriptsize$\pm$0.015}} \\
pima & 0.734{\scriptsize$\pm$0.053} & 0.653{\scriptsize$\pm$0.040} & 0.664{\scriptsize$\pm$0.040} & 0.673{\scriptsize$\pm$0.024} & 0.655{\scriptsize$\pm$0.021} & \textbf{0.675{\scriptsize$\pm$0.043}} \\
\bottomrule
\end{tabular}
}
\end{table}
